\chapter{Présentation du Groupe Managem}

Ce premier chapitre présente le Groupe Managem, entité dont dépend l'entreprise qui m'a accueillie. Les éléments qui suivent proviennent de sources publiques et de la documentation institutionnelle du Groupe~; ils fixent le décor dans lequel s'inscrit mon stage.

\section{Un acteur historique de la mine au Maroc}

Managem est un groupe minier marocain dont les origines remontent à la découverte d'un gisement de cobalt à Bou-Azzer, dans l'Anti-Atlas, à la fin des années 1920~; la première mine du Groupe démarre son activité au début des années 1930~\cite{managem_wiki}. Structuré sous sa forme actuelle dans les années 1990 pour regrouper les activités minières du groupe ONA, il est aujourd'hui coté à la Bourse de Casablanca et rattaché au holding Al~Mada~\cite{managem_wiki}.

Le Groupe couvre l'ensemble de la chaîne minière~: exploration, extraction, traitement (dont l'hydrométallurgie) et commercialisation de métaux de base, de métaux précieux, de cobalt et d'autres substances minérales~\cite{managem_wiki}. Il est présenté par plusieurs sources comme le leader du secteur minier métallique marocain~\cite{afd_potentiel}.

\section{Une présence panafricaine}

Au-delà du Maroc, Managem a développé depuis plus de vingt ans une stratégie d'expansion sur le continent. Le Groupe exploite un ensemble de mines réparties dans plusieurs pays d'Afrique --- notamment la Guinée, le Gabon, la République démocratique du Congo (RDC), le Soudan et, plus récemment, le Sénégal~\cite{afd_potentiel,ecofin_managem}. Cette dimension panafricaine est l'un des traits qui m'ont le plus frappée~: il ne s'agit pas d'une entreprise strictement nationale, mais d'un groupe qui pilote des projets sur des géographies très différentes.

\section{Des minerais stratégiques}

Le portefeuille du Groupe est centré sur plusieurs minerais à forte valeur, dont certains sont dits « stratégiques » car essentiels à la transition énergétique~:

\begin{itemize}
\item le \textbf{cobalt}, historiquement à l'origine du Groupe, avec le gisement de Bou-Azzer, l'un des rares au monde produisant principalement du cobalt~\cite{ecofin_managem}~;
\item l'\textbf{argent}, dont Managem est un producteur majeur en Afrique, notamment via le gisement d'Imiter~\cite{ecofin_managem}~;
\item le \textbf{cuivre}, au cœur de projets d'envergure comme Tizert~;
\item l'\textbf{or}, en forte croissance grâce aux projets ouest-africains (Guinée, Sénégal)~;
\item ainsi que le \textbf{zinc}, le \textbf{plomb} et des minéraux industriels comme la \textbf{fluorine}~\cite{morocco_minerals}.
\end{itemize}

\section{Des projets miniers d'envergure}

Parmi les grands chantiers récents du Groupe, on peut citer la mine d'or de Tri-K en Guinée, le projet cobalt-cuivre de Pumpi en RDC, la mine d'or de Boto au Sénégal et le projet cuprifère de Tizert au Maroc~\cite{ecofin_managem}. Ces deux derniers projets, dont j'ai entendu parler pendant mon stage, sont détaillés au chapitre~\ref{chap:projets}.

\section{Des filiales spécialisées}

Le Groupe fonctionne à travers un ensemble de filiales spécialisées, qui couvrent chacune un maillon de la chaîne. Parmi elles figurent des sociétés d'exploitation minière, une société métallurgique (production d'argent, notamment), des sociétés dédiées à l'or et au cuivre, et --- ce qui me concerne directement --- \textbf{Reminex}, la filiale d'ingénierie et de recherche du Groupe, présentée au chapitre suivant.

\begin{presente}
\textbf{Ce que je retiens de ce cadre.} Avant ce stage, je voyais « la mine » comme une simple activité d'extraction. J'ai compris que Managem est en réalité un groupe intégré, qui va de la recherche du gisement jusqu'à la vente du métal, sur plusieurs continents, et qui s'appuie pour cela sur des filiales aux métiers très différents.
\end{presente}
